\documentclass[11pt]{article}
\usepackage[margin=0.85in]{geometry}
\usepackage{amsmath,amssymb,bm}
\usepackage{xcolor}
\usepackage[most]{tcolorbox}
\usepackage{enumitem}
\usepackage{booktabs}
\usepackage{parskip}
\usepackage{hyperref}
\hypersetup{colorlinks=true,linkcolor=blue!50!black,urlcolor=blue!50!black}
\setlist{nosep,leftmargin=1.5em}

\newtcolorbox{keybox}[1][]{colback=blue!4,colframe=blue!45,boxrule=0.6pt,arc=2pt,
  left=7pt,right=7pt,top=5pt,bottom=5pt,#1}
\newtcolorbox{ideabox}[1][]{colback=green!5,colframe=green!50!black,boxrule=0.8pt,arc=2pt,
  left=7pt,right=7pt,top=6pt,bottom=6pt,#1}
\newtcolorbox{warnbox}[1][]{colback=red!4,colframe=red!50!black,boxrule=0.6pt,arc=2pt,
  left=7pt,right=7pt,top=5pt,bottom=5pt,#1}

\title{\vspace{-2.4em}\textbf{A New Sparse Array Geometry for\\ Near-Field Cell-Free XL-MIMO}\\[0.3em]
\large What the four group papers do, why our current array is not enough,\\ and four geometry ideas that could be a real contribution}
\author{}
\date{}

\begin{document}
\maketitle
\vspace{-3em}

\section{Purpose of this document}

Your professor's message was: the symmetric nested array is fine, but it is \emph{not a
contribution}. Symmetric co-arrays already exist. The contribution has to be in \textbf{how we
arrange the antennas} --- the array geometry itself. He pointed you at four papers from the group.

This document does four things, in simple English:
\begin{enumerate}
\item Explains what each of the four papers actually does.
\item Explains \emph{why} our current symmetric array performs the same as the standard difference
co-array (the result you observed) --- there is a concrete technical reason.
\item Explains what the near field really demands from an array geometry (the mathematics, step by
step, with every symbol defined).
\item Proposes \textbf{four new geometry ideas}, with an honest assessment of which one is the real
contribution and what the risks are.
\end{enumerate}

\section{What the four papers do}

\subsection{Paper 1 --- Pinto \& de Lamare, EUSIPCO 2023 (super-resolution arrays for multiuser MIMO)}

This is the paper your professor originally pointed you to. In plain words:
\begin{itemize}
\item Uplink of a \emph{single-cell} multiuser MIMO system. The base station uses a sparse array ---
either a two-level nested array (TLNA) or a coprime array (CPA).
\item They take the received covariance matrix $\mathbf{R}_x$, vectorize it
($\mathbf{v}=\mathrm{vec}(\mathbf{R}_x)$), remove duplicate entries, and sort them. This produces a
\emph{virtual} (difference co-array) array with far more sensors than the physical array.
\item They then apply \emph{spatial smoothing} to make the covariance full rank again, and run
\textbf{MMSE} and \textbf{ordered successive interference cancellation (OSIC)} receivers on the
enlarged virtual data.
\item \textbf{Headline result:} a sparse array with $M=8$ physical antennas matches or beats a
uniform array with $M=16$ antennas, in both achievable sum-rate and bit error rate. Half the
hardware, same performance.
\end{itemize}

\begin{warnbox}
\textbf{Important limitation for us:} this paper assumes a \emph{far-field}, single-path,
line-of-sight-style channel $\mathbf{h}_k=g_k\,\mathbf{a}(\theta_k)$, with no scattering and a
single cell. There is no near field and no distributed access points. So it cannot be used in our
scenario unchanged --- which is exactly the gap we are trying to fill.
\end{warnbox}

\subsection{Paper 2 --- Leite, de Lamare, Zakharov, Liu \& Haardt, 2024 (DOA with sparse subarrays)}

This is the most relevant paper to a \emph{cell-free} system, and I think it is the one the professor
most wants you to absorb.
\begin{itemize}
\item It studies arrays that are broken into \textbf{subarrays}. Two families are defined:
  \begin{itemize}
  \item \textbf{Type-I}: take one known sparse geometry and \emph{cut} it into pieces.
  \item \textbf{Type-II}: take one reference sparse subarray and \emph{repeat} it (translated copies)
  --- a ``multiple invariance'' structure.
  \end{itemize}
\item The key practical assumption is \textbf{partial calibration}: inside each subarray the sensors
are phase-synchronised, but the \emph{relative phase between subarrays is unknown}.
\item Algorithms \textbf{GCA-MUSIC} and \textbf{GCA-rMUSIC}: build the difference co-array of
\emph{each} subarray, spatially smooth it, get each subarray's noise subspace, and then combine them
by \emph{intersecting} the noise subspaces using projection onto convex sets (POCS). This is a
\emph{non-coherent} combination --- no phase synchronisation needed.
\item \textbf{Theorem 1} gives the total degrees of freedom of a type-II array:
\[
\mathrm{DoF}\;=\;L(\mathrm{sDoF}-1)+2(L-1)\mu+1 ,\qquad 1\le\mu\le\kappa,
\]
where $L$ is the number of subarrays, $\mathrm{sDoF}$ is the DoF of one subarray, $\mu$ is the
normalised spacing between subarrays, and $\kappa$ is the aperture of one subarray. Example in the
paper: $L=2$ subarrays of $7$ sensors ($N=14$ total) reach $\mathrm{DoF}=85$.
\item They can estimate \emph{more sources than sensors in each subarray}.
\end{itemize}

\begin{keybox}
\textbf{Why this matters enormously for us.} A cell-free network \emph{is} a sparse subarray system:
each access point (AP) is a subarray, and the APs are \emph{not} phase-synchronised with each other
--- exactly the ``partially calibrated'' case. Our CAPE technique already fuses APs non-coherently,
which is the same philosophy as GCA-MUSIC. So this paper is the natural theoretical home for our
work. The gap: Leite et al.\ assume the subarrays are \emph{collinear} and are \emph{translated
copies} of one another, and the whole treatment is \emph{far-field}. Our APs are scattered in two
dimensions and our users are in the near field.
\end{keybox}

\subsection{Papers 3 and 4 --- the MISC family (the array-design recipe)}

These two are the ones that show you \emph{how this group designs a new array geometry}. They are the
template you should copy.

\textbf{Paper 3 --- Shi, Li \& de Lamare, IEEE SPL 2022 (IMISC).}
\begin{itemize}
\item MISC = ``maximum inter-element spacing constraint''. The idea: do not describe the array by its
positions, describe it by its \textbf{inter-element spacing (IES) set} --- the list of gaps between
neighbouring antennas.
\item The IES set is fixed by just \emph{two} numbers: the number of sensors $Q$ and the maximum
allowed spacing $M$, with $M=4\lfloor (Q+2)/6\rfloor$.
\item The resulting IMISC array is a concatenation of \textbf{six uniform sub-ULAs}, each with its
own spacing.
\item They derive in closed form the \textbf{uniform degrees of freedom} (uDOF) --- the length of the
hole-free consecutive part of the difference co-array:
$\mathrm{uDOF}_{\mathrm{IMISC}}\approx \tfrac{2Q^2}{3}-\tfrac{2Q}{3}-1$, i.e.\ order $2Q^2/3$,
which beats the $4Q^2/7$ of the earlier arrays.
\item They also derive the \textbf{weight function} $w(n)$ = the number of antenna pairs separated by
exactly $n$ units. The first three values $w(1),w(2),w(3)$ dominate \textbf{mutual coupling} (MC).
For IMISC: $w(1)=2$, $w(2)=2\lfloor (Q+2)/6\rfloor$ or $5$, $w(3)=1$ or $2$.
\end{itemize}

\textbf{Paper 4 --- Sheng, Lu, Li \& de Lamare, IEEE TCAS-II 2024 (EMISC).}
\begin{itemize}
\item Same recipe, pushed further. EMISC uses \textbf{seven} uniform sub-ULAs instead of six (an
extra sub-ULA with spacing 3 is inserted and everything is re-optimised).
\item uDOF: $\mathrm{uDOF}_{\mathrm{EMISC}}=-\tfrac{3L^2}{2}+(2K-1)L+7$ with
$L=4\lfloor (K-4)/6\rfloor+4$ --- four more than IMISC.
\item The real trick: the rearrangement makes $\boxed{w(1)=1}$ --- only \emph{one} pair of adjacent
antennas in the whole array. Since $w(1)$ dominates mutual coupling, this gives the lowest coupling
of the family. Ordering: $\mathrm{CL}_{\mathrm{EMISC}}<\mathrm{CL}_{\mathrm{IMISC}}<\mathrm{CL}_{\mathrm{MISC}}$.
\end{itemize}

\begin{keybox}
\textbf{The recipe, extracted.} Every new-geometry paper in this group follows the same five steps:
\begin{enumerate}
\item Describe the array by an \textbf{IES set} (gaps), parameterised by a couple of integers.
\item Build the array as a \textbf{concatenation of uniform sub-ULAs}.
\item Prove which part of the difference co-array is \textbf{hole-free}, and get a
\textbf{closed-form uDOF}.
\item Derive the \textbf{weight function} $w(1),w(2),w(3)$ $\Rightarrow$ mutual coupling.
\item Show you beat the previous arrays on \emph{both} uDOF \emph{and} coupling.
\end{enumerate}
If we want a contribution the professor will accept, our new array must come with these same five
deliverables.
\end{keybox}

\section{Why our current symmetric array performs the same as the standard one}

You observed that the symmetric near-field sparse array gives almost identical BER, sum-rate and
NMSE to the standard difference co-array. That is not a coincidence or a bug. There are four
concrete reasons:

\begin{enumerate}
\item \textbf{Both use the same estimator.} In our code, both geometries do exactly the same thing:
difference-co-array MUSIC $\rightarrow$ angle estimate $\rightarrow$ a separate 1-D range search on
the physical array $\rightarrow$ feed the position to CAPE. Changing the geometry from ``nested'' to
``symmetric nested'' changes \emph{which lags exist}, but not \emph{what the algorithm extracts}. If
the processing does not exploit the symmetry, the symmetry buys nothing.

\item \textbf{The symmetric array is smaller.} At equal antenna count $N=24$, our symmetric nested
array has an aperture of $4.20$\,m while the plain nested array has $7.75$\,m. Angular resolution
scales with aperture, so we \emph{lost} resolution by symmetrising. Worse, the Rayleigh distance
$d_R = 2D^2/\lambda$ shrinks with aperture, so the symmetric array sees a \emph{smaller near-field
region} --- the opposite of what we want.

\item \textbf{It has fewer virtual sensors.} At $N=24$: plain nested gives 311 distinct lags, our
symmetric nested only 109. Symmetry \emph{wastes} degrees of freedom, because mirror pairs produce
\emph{duplicate} difference lags (each lag is counted twice), which raises the weight function
$w(n)$ and lowers uDOF.

\item \textbf{The co-array never produces the range.} The near-field part of the problem --- the
distance $r$ --- is recovered in our code by a one-dimensional search on the \emph{physical}
observation, not by the virtual array. So the co-array contributes nothing to the specifically
near-field part of the estimate.
\end{enumerate}

\begin{warnbox}
\textbf{Conclusion.} ``Make it symmetric'' is a \emph{validity} fix, not a \emph{performance} fix. It
stops the difference co-array from being mathematically wrong in the near field, but it gives up
aperture and degrees of freedom to do so, and the two effects cancel. To get a real gain, the
geometry must be designed so that the \emph{near-field information itself} (angle \emph{and} range)
is carried by the virtual array --- and that is a genuine, unsolved design problem.
\end{warnbox}

\section{What the near field actually demands from a geometry}

\subsection{The near-field phase, step by step}

Let an antenna sit at position $x_m$ on a line (metres), and a user sit at distance $r$ (metres) and
angle $\theta$. The exact distance from user to antenna $m$ is
\begin{equation}
r_m=\sqrt{r^2+x_m^2-2 r x_m\sin\theta}.
\end{equation}
Expanding this for $r$ larger than $x_m$ but not enormously larger (the Fresnel / near-field region):
\begin{equation}
r_m\;\approx\; r-x_m\sin\theta+\frac{x_m^2\cos^2\theta}{2r}.
\end{equation}
\emph{Here $r$ is the user distance, $\theta$ the angle, $x_m$ the antenna position.} The received
phase at antenna $m$ is $-\tfrac{2\pi}{\lambda}r_m$, so dropping the common term $-\tfrac{2\pi}{\lambda}r$
the array response is
\begin{equation}
a_m=\exp\!\big(j(\omega x_m+\psi x_m^2)\big),\qquad
\omega \triangleq \frac{2\pi}{\lambda}\sin\theta,\qquad
\psi \triangleq -\frac{\pi\cos^2\theta}{\lambda\, r}.
\end{equation}
\emph{Here $\lambda$ is the wavelength, $\omega$ is the ``angle parameter'' (linear in position) and
$\psi$ is the ``range parameter'' (multiplies position \emph{squared}). In the far field $r\to\infty$
so $\psi\to 0$ and only the linear term survives --- that is the classical case.}

\subsection{Why the difference co-array breaks, in one equation}

The covariance between antennas $m$ and $n$ for one user of power $\sigma^2$ is
\begin{equation}
[\mathbf{R}]_{m,n}=\sigma^2 a_m a_n^{*}
=\sigma^2\exp\!\Big(j\big[\underbrace{\omega(x_m-x_n)}_{\text{difference}}
+\psi\underbrace{(x_m^2-x_n^2)}_{=(x_m-x_n)(x_m+x_n)}\big]\Big).
\end{equation}
Write $d=x_m-x_n$ (the \textbf{difference lag}) and $s=x_m+x_n$ (the \textbf{sum lag}). Then
\begin{equation}
\boxed{\;[\mathbf{R}]_{m,n}=\sigma^2\exp\!\big(j[\,\omega d+\psi\, d\, s\,]\big).\;}
\end{equation}
\emph{Read this carefully, it is the whole story:} the entry depends on the difference lag $d$
\textbf{and} on the product $d\,s$. Two antenna pairs with the \emph{same} difference $d$ but
\emph{different} sums $s$ give \emph{different} phases. That is precisely why you cannot merge them
into one ``virtual sensor'' --- the classical difference co-array assumption fails.

\subsection{The two escapes, and the tension between them}

\textbf{Escape 1 --- symmetric pairs give clean angle.} If the array contains both $+x$ and $-x$,
then the pair $(m,n)=(x,-x)$ has $s=0$, so the quadratic term vanishes:
\begin{equation}
[\mathbf{R}]_{x,-x}=\sigma^2\exp(j\,\omega\,2x).
\end{equation}
This is a pure far-field-like term with difference lag $2x$. \emph{This is the real reason symmetric
arrays are used in the near field}: the mirror pairs give an angle-only virtual array, immune to
range curvature.

\textbf{The catch:} the set of clean angle lags is $\{2x : x \text{ and } -x \text{ both present}\}$,
i.e.\ \emph{twice the positive positions}. For this set to be hole-free (consecutive), the positive
positions must themselves be consecutive on a fine grid --- which means a \emph{dense}, small array.
This is exactly why our symmetric nested array came out compact with a small aperture. \textbf{The
requirement ``symmetric and hole-free'' fights the requirement ``large aperture''.}

\textbf{Escape 2 --- the sum co-array carries the range.} Pairs with the same $d$ but different $s$
differ in phase by $\psi d (s_1-s_2)$. That difference is \emph{pure range information}. So the
\textbf{sum co-array} $\mathbb{S}=\{x_m+x_n\}$ is not a nuisance --- it is the resource that lets us
estimate $r$. A good near-field geometry should make \emph{both} the difference co-array (for angle)
and the sum co-array (for range) rich and hole-free.

\subsection{The design requirements, collected}

\begin{keybox}
A sparse array for \textbf{near-field cell-free XL-MIMO} must satisfy, simultaneously:
\begin{description}[leftmargin=3.2em,style=nextline]
\item[\textbf{R1 Aperture (new!)}] $D\ \ge\ \sqrt{\lambda\, r_{\max}/2}$, so that users at distance
up to $r_{\max}$ are genuinely inside the near field ($r<2D^2/\lambda$). \emph{No MISC-family paper
has this constraint} --- they only cap the \emph{maximum} spacing. This is a two-sided problem:
aperture big enough for near field, spacing small enough for coupling and ambiguity.
\item[\textbf{R2 Angle DOF}] A long hole-free set of \emph{clean} angle lags (symmetric pairs,
$s=0$).
\item[\textbf{R3 Range DOF}] A rich, hole-free \emph{sum} co-array $\{x_m+x_n\}$, so range is
observable from the virtual array rather than from a separate search.
\item[\textbf{R4 Low coupling}] Small $w(1),w(2),w(3)$ --- ideally $w(1)=1$, the EMISC trick.
\item[\textbf{R5 Ambiguity}] Grating lobes controlled --- every sparse array creates them.
\item[\textbf{R6 Distributed}] Across APs: partially calibrated (no phase sync), so fusion must be
non-coherent, and the per-AP geometries should \emph{complement} each other.
\end{description}
\end{keybox}

\section{Four geometry ideas}

\subsection{Idea 1 (recommended) --- the Two-Scale Symmetric-Core MISC array (TS-MISC)}

\begin{ideabox}
\textbf{The structure.} Build the array from three functional parts instead of one:
\begin{enumerate}
\item A \textbf{small symmetric dense core} around the origin, on a \emph{quarter-wavelength}
($\lambda/4$) grid: positions $\pm\{\tfrac{1}{2},1,\tfrac{3}{2},\dots\}$ in units of
$d=\lambda/2$. Its mirror pairs produce a \emph{consecutive} set of clean angle lags (R2). It is
cheap: only a handful of elements.
\item A \textbf{sparse asymmetric outer pair of wings}, placed by the \textbf{EMISC IES rule}
(sub-ULAs with prescribed gaps). The wings deliver the \textbf{aperture} (R1), the high
difference-co-array uDOF, and they are placed so that $w(1)=1$ overall (R4).
\item The \textbf{cross pairs} between core and wings have large, \emph{distinct} sum lags $s$, which
is exactly what makes the \textbf{range} observable from the virtual array (R3).
\end{enumerate}
\textbf{Why it fixes our problem.} It separates the two fighting requirements instead of compromising
between them: the \emph{core} is symmetric and dense (so the angle co-array is hole-free), the
\emph{wings} are sparse and wide (so the aperture and near-field reach are large). Our current
array tries to be symmetric \emph{everywhere}, which is why it ended up small and low-DOF.
\end{ideabox}

\textbf{What you would have to deliver} (the five MISC steps): a closed-form IES set parameterised by
$(K, L, \text{core size})$; the position set; a proof of the hole-free part of the \emph{difference}
co-array giving $\mathrm{uDOF}_{\mathrm{diff}}$; a proof of the hole-free part of the \emph{sum}
co-array giving $\mathrm{uDOF}_{\mathrm{sum}}$ (this is the genuinely new quantity); and the weight
function $w(1),w(2),w(3)$.

\subsection{Idea 2 --- Distributed Co-Array Tiling across APs (the cell-free contribution)}

\begin{ideabox}
\textbf{The structure.} Stop designing \emph{one} array. Design the \textbf{set of $L$ AP arrays
jointly}, so that the lags one AP is missing are supplied by another.
\begin{itemize}
\item Give AP $\ell$ its own IES set $\mathbb{S}_\ell$, chosen so that the \emph{union} of the
per-AP difference co-arrays, $\bigcup_\ell \mathbb{D}_\ell$, is hole-free and long --- while each
individual AP stays small and cheap.
\item Because APs are not phase-synchronised, the combination must be \emph{non-coherent}: exactly
the GCA-MUSIC / POCS machinery of Paper 2, and exactly what our CAPE already does.
\item This generalises Leite et al.'s type-II arrays, where every subarray is an identical
\emph{translated copy}. Here the subarrays are deliberately \emph{different} and
\emph{complementary}, and they are not collinear.
\end{itemize}
\textbf{Why it is a strong contribution.} It is the first array-geometry design where the unit of
design is the \emph{network}, not the panel. It fits our system perfectly and it is a natural
``type-III'' extension of the group's own framework.
\end{ideabox}

\subsection{Idea 3 --- RETRACTED, and the lemma that replaces it}

I originally proposed a ``half-shifted interleaved symmetric array'': put half the elements on the
integer grid and half on a half-integer grid, so that the clean angle lags $2x$ would fill both odd
and even values. \textbf{I tested this numerically and it is wrong}, so I am retracting it rather
than leaving it in. The reason is a simple scaling argument: if every position is an integer
multiple of \emph{some} unit $u$, then $2x$ is an \emph{even} multiple of $u$ no matter what $u$ is.
Shifting the grid by half a unit just renames the unit; it creates no new lags. Measured on a
$\lambda/4$ grid, the ``clean-angle'' run length stayed at $1$ for every array I tried, including
the half-shifted ones.

What survives is a cleaner and more useful statement:

\begin{keybox}
\textbf{Lemma (mirror-pair bottleneck).} The clean angle co-array --- the set of lags obtained from
mirror pairs, $\{\,2x : x\in\mathbb{S},\ -x\in\mathbb{S}\,\}$ --- is hole-free \textbf{if and only
if the positive half of the array is a contiguous ULA}.

\emph{Proof sketch:} the clean lag set is exactly $2\mathbb{S}^{+}$ where $\mathbb{S}^{+}$ is the
positive half; $2\mathbb{S}^{+}$ is consecutive on the $2u$ lattice iff $\mathbb{S}^{+}$ is
consecutive on the $u$ lattice. $\square$
\end{keybox}

\textbf{Why this lemma matters.} It says the tension in Section 4.3 is not an artefact of a bad
design --- it is a \emph{structural law}. You cannot have a hole-free clean-angle co-array without a
dense symmetric core. Any near-field sparse array must therefore \emph{spend} some elements on a
dense core and buy aperture with the rest. That is the real justification for the core-plus-wings
structure of Idea 1, and stating and proving this lemma is itself a small theoretical contribution.

\subsection{Idea 4 --- Curvature-Matched Multi-Scale (geometric-progression) array}

\begin{ideabox}
\textbf{The observation.} Range information lives in the \emph{curvature} of the phase across the
array --- the $\psi x^2$ term. Curvature is best measured when you have element separations at
\emph{many different scales}, not one.

\textbf{The structure.} Choose the IES set as a \emph{geometric} progression,
$\{1,2,4,8,\dots\}\times d$ (possibly mirrored), rather than the arithmetic progressions used by
nested, coprime and MISC arrays. This samples the quadratic phase term at many scales with very few
elements, making the range parameter $\psi$ well-conditioned.
\end{ideabox}

\textbf{Honest caveat.} Geometric/fractal spacings give an excellent spread of \emph{sum} lags but a
\emph{hole-riddled} difference co-array, so angle DOF suffers. This idea is best used as the
\emph{wing} rule inside Idea 1, not on its own.

\section{I tested Idea 1 numerically --- here is the honest result}

Before asking you to invest in this, I built a prototype TS-MISC array and measured it against the
alternatives at \textbf{equal antenna count $N=24$}. All lengths use $\lambda=0.1$\,m. ``uDOF'' is
the longest hole-free run of the co-array on the $\lambda/2$ lattice; $w(1)$ counts adjacent pairs
(mutual coupling); $d_R=2D^2/\lambda$ is the Rayleigh distance (near-field reach).

\begin{center}
\begin{tabular}{@{}lrrrrrrc@{}}
\toprule
\textbf{Array} & \textbf{$N$} & \textbf{$D$ [m]} & \textbf{$d_R$ [m]} & \textbf{uDOF$_{\rm diff}$} &
\textbf{uDOF$_{\rm sum}$} & \textbf{$w(1)$} & \textbf{sym.?} \\
\midrule
Dense ULA               & 24 & 1.15 & 26   & 47  & 47  & 23 & yes \\
Nested (standard)       & 24 & 7.75 & 1201 & \textbf{311} & 168 & 12 & no \\
Symmetric nested (ours) & 24 & 4.20 & 353  & 99  & 99  & 12 & yes \\
TS-MISC prototype       & 24 & 7.60 & 1155 & 69  & 69  & \textbf{6} & core \\
\bottomrule
\end{tabular}
\end{center}

\begin{warnbox}
\textbf{Read this honestly: my prototype does not yet win.} On the headline metric
(uDOF$_{\rm diff}$) the plain nested array is far ahead --- $311$ versus my $69$. I am not going to
present the idea as proven when the numbers say otherwise.

What the table \emph{does} establish is worth having:
\begin{itemize}
\item \textbf{Our current symmetric nested array is confirmed as a bad trade.} Symmetrising cost us
$311\rightarrow 99$ in uDOF \emph{and} $7.75\rightarrow 4.20$\,m of aperture. That is the measured
explanation for why it performs the same as the standard co-array --- it gave up two things to gain
one.
\item \textbf{The plain nested array has \emph{zero} clean-angle lags} (it is not symmetric at all),
so despite its uDOF of 311 it is not formally valid in the near field. The two baselines each fail a
different requirement.
\item \textbf{The core-plus-wings structure does deliver on two of the four axes:} it keeps nearly
the full aperture ($7.60$ vs $7.75$\,m, so near-field reach is preserved) and it \textbf{halves the
mutual coupling} ($w(1)=6$ versus $12$) while retaining a symmetric core.
\end{itemize}
\textbf{The open problem --- and this is exactly where the contribution lies --- is the wing
placement rule.} I placed the wings by hand; a real design must derive them from an EMISC-style IES
set so that uDOF climbs back toward (and past) $311$ while $w(1)$ stays at $6$ or below and the core
stays symmetric. \emph{That derivation is the research work, and it is not yet done.}
\end{warnbox}

\section{Which one to choose, and the honest risks}

\begin{center}
\begin{tabular}{@{}p{3.1cm}p{3.4cm}p{3.4cm}p{3.6cm}@{}}
\toprule
\textbf{Idea} & \textbf{Main strength} & \textbf{Main weakness} & \textbf{Novelty risk} \\
\midrule
1. TS-MISC core+wings & Fixes the exact problem we measured; follows the group's recipe & Needs two
new proofs (diff \emph{and} sum uDOF) & Medium: near-field sparse arrays exist, but not via
IES/MISC construction \\
2. Distributed tiling & Unique to cell-free; matches CAPE and Paper 2 & Harder to analyse; needs a
network-level DOF theorem & \textbf{Low risk: genuinely unexplored} \\
3. \emph{(retracted)} & Replaced by the mirror-pair lemma & --- & Idea was wrong; lemma stands \\
4. Multi-scale wings & Best range conditioning & Poor angle DOF alone & Medium: fractal arrays exist
(far-field) \\
\bottomrule
\end{tabular}
\end{center}

\begin{keybox}
\textbf{Recommendation.} Make \textbf{Idea 1 the array contribution} and \textbf{Idea 2 the system
contribution}, and use Idea 3/4 as design rules \emph{inside} Idea 1. That gives a paper with two
linked novelties:
\begin{itemize}
\item a new \emph{single-AP} geometry (TS-MISC) that is the first MISC-style array designed for the
near field, with a closed-form IES set, a \emph{joint} difference-and-sum uDOF analysis, and a
minimum-aperture criterion that no existing MISC paper has;
\item a new \emph{network-level} design (distributed co-array tiling) where the APs' geometries are
complementary and fused non-coherently --- a true ``type-III'' sparse subarray.
\end{itemize}
\end{keybox}

\begin{warnbox}
\textbf{Novelty must be verified, not assumed.} My literature check found: (i) near-field sparse
arrays \emph{do} exist --- in particular a design based on symmetric augmentation of three ULAs for
near-field DOA; (ii) sum--difference co-array designs exist but are \emph{far-field}; (iii) I found
\emph{no} work combining the near-field spherical model with a joint hole-free sum--difference
design, and \emph{no} work designing co-array geometry across distributed/cell-free APs. Before
claiming novelty, search specifically for: ``near-field sum-difference coarray array design'',
``symmetric nested array near-field DOF'', and ``sparse array design cell-free''.
\end{warnbox}

\section{How we would prove it works, in our code}

The array contribution has to be measured with array-level metrics first, and only then with
communication metrics:
\begin{enumerate}
\item \textbf{Geometry metrics} (no simulation needed): $\mathrm{uDOF}_{\mathrm{diff}}$,
$\mathrm{uDOF}_{\mathrm{sum}}$, aperture $D$, Rayleigh distance $2D^2/\lambda$, near-field link
fraction, and the weight function $w(1),w(2),w(3)$ (coupling). Compare against ULA, nested, coprime,
ECA, MISC, IMISC, EMISC at \emph{equal antenna count}.
\item \textbf{Estimation metrics}: angle RMSE, range RMSE, channel NMSE, CAPE localisation error.
\item \textbf{Communication metrics}: BER and sum rate for the four detectors, fronthaul and
complexity.
\end{enumerate}

\textbf{Two code issues to fix at the same time} (you asked why they happen):
\begin{itemize}
\item \emph{Symmetric array ties the standard co-array} --- explained in Section 3: same estimator,
smaller aperture, fewer lags, and range never comes from the co-array. Fixing it requires the new
geometry \emph{and} an estimator that uses the sum co-array for range.
\item \emph{Cross-AP list detection equals plain list detection under estimated CSI} --- the most
likely cause is that CAPE has already consumed the cross-AP diversity (all APs are handed the same
fused position, so their channel estimates are correlated), leaving the cross-AP step no independent
information to exploit; a secondary cause is the list pruning threshold collapsing the list to a
single survivor. Both are testable: print the average surviving list size, and compare cross-AP gain
with CAPE on versus off.
\end{itemize}

\section{Suggested next step}

\begin{enumerate}
\item Fix the TS-MISC construction on paper: choose the core size, the wing IES rule, write the
position set in closed form.
\item Compute $\mathrm{uDOF}_{\mathrm{diff}}$, $\mathrm{uDOF}_{\mathrm{sum}}$ and $w(1),w(2),w(3)$
numerically for $K=10\dots60$ and compare with nested / coprime / ECA / IMISC / EMISC.
\item Only if the geometry numbers win, put it into the simulator and measure NMSE, BER, sum rate.
\item In parallel, formulate the network-level (Idea 2) DOF statement, mirroring Theorem 1 of
Paper 2.
\end{enumerate}

\begin{thebibliography}{9}
\bibitem{pinto2023} S. F. B. Pinto and R. C. de Lamare, ``Multiuser multiple-antenna wireless
communications systems based on super-resolution arrays,'' \emph{EUSIPCO}, 2023, pp. 1464--1468.
\bibitem{leite2024} W. S. Leite, R. C. de Lamare, Y. Zakharov, W. Liu and M. Haardt, ``Direction of
arrival estimation with sparse subarrays,'' arXiv:2409.00033, 2024.
\bibitem{shi2022} W. Shi, Y. Li and R. C. de Lamare, ``Novel sparse array design based on the maximum
inter-element spacing criterion,'' \emph{IEEE Signal Process. Lett.}, vol. 29, pp. 1754--1758, 2022.
\bibitem{sheng2024} X. Sheng, D. Lu, Y. Li and R. C. de Lamare, ``Enhanced MISC-based sparse array
with high uDOFs and low mutual coupling,'' \emph{IEEE Trans. Circuits Syst. II}, vol. 71, no. 2,
pp. 972--976, 2024.
\bibitem{pal2010} P. Pal and P. P. Vaidyanathan, ``Nested arrays: a novel approach to array
processing with enhanced degrees of freedom,'' \emph{IEEE Trans. Signal Process.}, vol. 58, no. 8,
pp. 4167--4181, 2010.
\bibitem{zheng2019} Z. Zheng, W. Q. Wang, Y. Kong and Y. D. Zhang, ``MISC array: a new sparse array
design achieving increased degrees of freedom and reduced mutual coupling effect,'' \emph{IEEE Trans.
Signal Process.}, vol. 67, no. 7, pp. 1728--1741, 2019.
\bibitem{liu2016} C. L. Liu and P. P. Vaidyanathan, ``Super nested arrays: linear sparse arrays with
reduced mutual coupling --- Part I,'' \emph{IEEE Trans. Signal Process.}, vol. 64, no. 15,
pp. 3997--4012, 2016.
\bibitem{nfsparse} ``Near-field DOA estimation oriented sparse array design via symmetric
augmentation of three ULAs,'' \emph{Digital Signal Processing}, 2024.
\bibitem{sparseisac} ``Sparse MIMO for ISAC: new opportunities and challenges,'' arXiv:2406.12270,
2024.
\end{thebibliography}

\end{document}
